% =====================================================================
%  Capitolo 4 completo: Fisher-Kolmogorov sul connettoma cerebrale.
%  Si ferma subito prima del capitolo HPC su reordering e parallelismo.
%
%  COME USARE QUESTO FILE
%  Sopra ogni sottosezione c'e' un blocco "DA LEGGERE" con le pagine
%  esatte dei tre paper di riferimento, e sopra ogni figura un blocco
%  "ANALOGA" che indica la figura corrispondente nella letteratura.
%  Sono commenti LaTeX: non compaiono nel PDF compilato.
%
%  RIFERIMENTI, con le abbreviazioni usate nei blocchi:
%    [F]  Fornari, Schafer, Jucker, Goriely, Kuhl,
%         "Prion-like spreading of Alzheimer's disease within the brain's
%         connectome", J. R. Soc. Interface 16 (2019) 20190356, 12 pp.
%         Le pagine citate sono le pagine dell'articolo, che coincidono
%         con le pagine del PDF in QuantumGraphs/fornari.pdf.
%    [W]  Weickenmeier, Jucker, Goriely, Kuhl, "A physics-based model
%         explains the prion-like features of neurodegeneration...",
%         J. Mech. Phys. Solids 124 (2019) 264-281.
%         Cito le pagine di rivista; nel PDF in QuantumGraphs/Weickenmeier.pdf
%         la pagina di rivista p corrisponde alla pagina PDF p-263.
%    [C]  Corti, Antonietti, Bonizzoni, Dede, Quarteroni, Fisher-Kolmogorov
%         stocastico su connettoma. Le pagine citate sono quelle del PDF in
%         QuantumGraphs/corti.pdf.
%
%  Le chiavi \cite{} sono segnaposto: fornari2019, weickenmeier2019,
%  corti2023. Sostituirle con quelle della bibliografia del report, che
%  gia' cita Weickenmeier come [7].
%
%  Ogni numero riportato qui e' stato ricalcolato dai file di risultato
%  con scripts/verify-figures.py (91 controlli, nessuno fallito).
% =====================================================================


% ---------------------------------------------------------------------
%  DA LEGGERE prima dell'introduzione del capitolo
%    [F] par. 1 "Motivation", pp. 1-2: perche' un modello di rete invece
%        di un continuo, e il confronto di costo con gli elementi finiti
%        tridimensionali (400 000 tetraedri, 80 000 gradi di liberta').
%    [W] par. 1 "Motivation", pp. 265-266: il paradigma prionico e le tre
%        malattie neurodegenerative.
%    [C] par. 1 "Introduction", pp. 1-2: la motivazione clinica e il ruolo
%        della quantificazione dell'incertezza.
%  Non serve leggere altro per l'introduzione: le tre motivazioni sono
%  sovrapponibili e il nostro paragrafo introduttivo e' gia' scritto.
% ---------------------------------------------------------------------

\section{Prion-like spreading of neurodegenerative diseases within the
brain's connectome}

We have validated our finite element library for the solution of graph-based
differential equations on a sequence of elliptic and parabolic benchmark
problems. We now consider a more application-oriented test case. In particular,
we study the Fisher--Kolmogorov equation on a brain connectome's topology. This
problem provides a natural extension of the previous linear parabolic models,
as it combines diffusion along the graph with a nonlinear reaction term. At the
same time, it provides a meaningful benchmark for assessing the scalability of
our framework through graph reordering and parallel programming techniques,
with the particularly interesting perspective of modeling the propagation of
infectious misfolded proteins through the brain, a process that is responsible
for neurodegenerative diseases. The goal of this section is therefore not to
provide a detailed biological study, but to investigate how the developed
solver can be extended from controlled benchmark problems to a realistic
network-based diffusion--reaction model.

The three works we follow play distinct roles, and it is worth separating them
at the outset. Weickenmeier et al.~\cite{weickenmeier2019} provide the
one-dimensional sensitivity test with which we validate the nonlinear solver
before it touches a connectome. Fornari et al.~\cite{fornari2019} provide the
83-region topology, the connectivity-weighted graph Laplacian, the entorhinal
seed, the parameters and the biomarker comparison. Corti et
al.~\cite{corti2023} provide the deterministic model with region-dependent
reaction coefficients and the semi-implicit time discretization we adopt. The
P1 metric-graph spatial discretization is the contribution of this work.


% ---------------------------------------------------------------------
%  DA LEGGERE prima di 4.1
%    [W] par. 2.2 "Continuum model of prion-like spreading", p. 267:
%        equazione (8), che e' la Fisher-Kolmogorov che risolviamo.
%    [W] par. 2.3 "Computational model", p. 268: residuo (10) e matrice
%        tangente (11). E' esattamente il nostro Newton su Eulero
%        implicito; vale la pena confrontare (11) con il nostro
%        assemblaggio del Jacobiano.
%    [W] par. 2.4 "Illustration of prion-like spreading", p. 268:
%        i parametri del test monodimensionale, nel = 200 elementi
%        lineari, T = 20, nstep = 200, alpha in {1,2,4}, d in {1,2,4}e-4,
%        c0 = 0.1 nel nodo centrale. E' il paragrafo da cui abbiamo preso
%        tutti i parametri, riga per riga.
%        Nello stesso paragrafo, la frase sullo studio preliminare del
%        passo temporale ("we had systematically varied the time step size
%        between dt/4 and 4 dt...") e' l'unico riferimento allo studio che
%        riproduciamo in Figura~\ref{fig:timestep}: nel paper non esiste
%        una figura corrispondente.
% ---------------------------------------------------------------------

\subsection{One-dimensional sensitivity analysis}
\label{sec:onedim}

Prion diseases are characterized by a chain reaction in which misfolded
proteins force native ones into a similar pathogenic structure. A physics-based
reaction--diffusion model can explain the growth and spreading of misfolded
protein in neurodegenerative diseases, like Alzheimer's disease. To
characterize the progression of misfolded proteins across the brain, we combine
the classical Fisher--Kolmogorov equation for population dynamics with an
anisotropic diffusion model. Following the work described
in~\cite{weickenmeier2019}, we perform a systematic sensitivity analysis,
demonstrating the role of the individual model parameters. We show how the
misfolded protein concentration is sensitive to the coefficients of growth and
extracellular diffusion. We solve
\begin{equation}
\frac{\partial c}{\partial t} - d\,\frac{\partial^2 c}{\partial x^2}
  = \alpha\,c\,(1-c),
\qquad (x,t) \in (-1,1) \times (0,20],
\label{eq:onedim}
\end{equation}
with homogeneous Neumann boundary conditions,
\[
\frac{\partial c}{\partial x}(-1,t) = \frac{\partial c}{\partial x}(1,t) = 0 .
\]
The initial concentration is zero at every mesh node except at the central
node, where $c(0,0)=0.1$. The interval is discretized using 200 linear finite
elements. Backward Euler is employed with $\Delta t = 0.1$, and the nonlinear
problem arising at each time step is solved by Newton's method. The sensitivity
analysis combines the diffusion coefficients
$d \in \{10^{-4}, 2\cdot10^{-4}, 4\cdot10^{-4}\}$ with the growth coefficients
$\alpha \in \{1,2,4\}$. Every one of these choices is the one made
in~\cite{weickenmeier2019}, including the number of elements and the number of
time steps.

% ---------------------------------------------------------------------
%  ANALOGA di Figura~\ref{fig:sensitivity}
%    [W] Fig. 3, p. 269. Stessa griglia 3x3, stesso significato dei
%        colori: per ogni coppia (x,c) il colore e' il primo istante in
%        cui la concentrazione in x raggiunge il livello c. Il nostro
%        pannello e' costruito con la stessa convenzione, non con una
%        rappresentazione diversa.
%    Da confrontare in particolare: la simmetria rispetto a x=0 e il
%    fatto che aumentando alpha da sinistra a destra i contorni si
%    addensano, mentre aumentando d dall'alto in basso si allargano.
% ---------------------------------------------------------------------

Figure~\ref{fig:sensitivity} presents the results using an activation-time
representation. For every pair $(x,c)$, the color denotes the first time at
which the concentration at position $x$ reaches the prescribed level $c$. The
black curves are therefore contours of equal activation time.

Increasing $\alpha$ accelerates the local conversion of healthy proteins,
whereas increasing $d$ enhances their spatial propagation. Consequently, both
parameters increase the propagation speed, although through different
mechanisms. The numerical fronts remain symmetric with respect to $x=0$: over
the whole dataset the largest left--right asymmetry is $3.33\cdot10^{-16}$, at
the level of round-off. The concentration range over all nine runs is
$[-6.34\cdot10^{-10}, 1.00]$, the negative value being numerical noise of the
same order.

\begin{figure}[htbp]
\centering
\includegraphics[width=\textwidth]{images/sensitivity.pdf}
\caption{Activation-time maps for the one-dimensional Fisher--Kolmogorov
problem. The growth coefficient increases from left to right, while the
diffusion coefficient increases from top to bottom. For each pair $(x,c)$, the
color represents the first time at which the numerical concentration at $x$
reaches the level $c$; black curves denote equal activation times. All panels
use the same time scale $t \in [0,20]$.}
\label{fig:sensitivity}
\end{figure}

% ---------------------------------------------------------------------
%  NOTA sulla velocita' del fronte: figura RIMOSSA, e il motivo e' tecnico.
%  Avevamo una figura a due pannelli che confrontava la velocita' misurata
%  con l'asintotico di Fisher-KPP. E' stata tolta perche' non e' una
%  verifica: l'eccesso rispetto a 2 sqrt(d alpha) vale il 2-6%, cresce con
%  d, dipende dalla finestra di fit per oltre un punto percentuale e NON
%  diminuisce restringendo il fit ai tempi tardi. Non si puo' nemmeno
%  mostrare che converga, perche' raffinando Delta t la discretizzazione
%  peggiora prima che la verifica converga (vedi Sezione 4.5.1).
%  Il dato quantitativo resta in benchmarks/18.../README.md come
%  diagnostica. Qui restiamo sul confronto qualitativo con [W], che e'
%  esattamente il livello a cui il paper stesso si ferma.
%  DA LEGGERE: [W] par. 2.4, p. 268, ultime righe, dove le tre proprieta'
%  qualitative del fronte sono elencate.
% ---------------------------------------------------------------------

The fronts remain symmetric about $x=0$ and their speed increases with both
$d$ and $\alpha$, consistently with the Fisher--KPP relation
$v_{\mathrm{front}} = 2\sqrt{d\alpha}$ and with the qualitative description
in~\cite{weickenmeier2019}. We also measured the front speed and found it to
exceed the asymptotic value by between two and six percent, growing with the
diffusivity, but we do not present that measurement as a verification: the
deviation does not decrease when the fit is restricted to later times, and the
time step cannot be refined far enough to settle the question, for the reason
given in Section~\ref{sec:boundary}. The quantitative diagnostic is retained in
the benchmark record rather than reported here.

% ---------------------------------------------------------------------
%  DA LEGGERE prima di questo paragrafo
%    [W] par. 2.4, p. 268, la frase: "In a preliminary convergence study, we
%        had systematically varied the time step size between [dt/4, dt/2, dt,
%        2 dt, 3 dt, 4 dt] with dt = 0.1, for a growth of alpha = 2 and a
%        spreading of d = 0.0002, and confirmed that the solution had converged
%        for dt = 0.1, while we observed a spurious increase in spreading for
%        larger time step sizes."
%        E' l'unica fonte. Il paper NON mostra ne' dati ne' figura, quindi non
%        possiamo ricostruire con quale criterio o quale norma abbiano deciso
%        che dt = 0.1 fosse convergente. Il testo qui sotto lo dice
%        esplicitamente, e per questo non afferma mai che la loro conclusione
%        sia sbagliata.
%
%  ATTENZIONE, distinzione da tenere ferma in tutta la sottosezione:
%    questo e' uno studio di SENSIBILITA' al passo temporale, non uno studio di
%    convergenza. La run piu' fine e' un riferimento numerico, non una
%    soluzione esatta. Gli ordini temporali formali sono verificati subito
%    dopo, contro la soluzione logistica esatta: primo ordine per Eulero
%    implicito, secondo per il semi-implicito di Crank-Nicolson.
%
%  ANALOGA di Figura~\ref{fig:timestep}
%    NESSUNA figura in [W]. Il pannello (a) e' la controparte grafica
%    dell'affermazione verbale sullo "spurious increase in spreading"; il
%    pannello (b) e' interamente nostro e non ha corrispondente.
% ---------------------------------------------------------------------

\subsubsection{Time-step sensitivity}
\label{sec:timestep}

The reference reports a preliminary study in which the time step was varied
between $\Delta t/4$ and $4\Delta t$ around $\Delta t = 0.1$, at $\alpha = 2$
and $d = 2\cdot10^{-4}$, concluding that $\Delta t = 0.1$ was regarded as
sufficiently converged and that larger steps produce a spurious increase in
spreading. That study is described in words only: neither the data nor a figure
is given, so the criterion and the norm behind the statement cannot be
reconstructed, and nothing in what follows should be read as a claim that the
statement is incorrect. What we can do is repeat the experiment with a stated
diagnostic and report what it shows.

We adopt as diagnostic the position $x_{\mathrm{f}}$ of the right front,
defined as the point where the profile crosses $c = 0.5$. All six runs are
compared at the common time $T = 19.2$, the largest time that is a whole
multiple of every step tested, including $\Delta t = 0.3$.

Figure~\ref{fig:timestep} shows the outcome, and it separates into two
statements of very different strength.

The qualitative observation of the reference is confirmed without reservation,
and it is the main conclusion of this subsection. Increasing the time step
moves the front outward, and beyond a point the front disappears altogether:
for $\Delta t \ge 0.2$ the interval is essentially full, with a minimum
concentration of $0.9320$ at $\Delta t = 0.2$ and of $1$ to four decimals at
$\Delta t = 0.3$ and $0.4$, so no $c = 0.5$ crossing exists at all. Large time
steps do produce a spurious increase in numerical propagation, exactly as
described.

The quantitative reading is more delicate. With the front position as
diagnostic, the three steps that still admit a front give
\[
x_{\mathrm{f}} = 0.6889, \quad 0.7268, \quad 0.8145
\qquad\text{for}\qquad
\Delta t = 0.025,\ 0.05,\ 0.1,
\]
so the front at $\Delta t = 0.1$ sits about $18.2\%$ farther out than at
$\Delta t = 0.025$. Under this diagnostic the solution at $\Delta t = 0.1$ is
therefore not yet insensitive to further temporal refinement.

This is not a contradiction of the reference, and it is worth being precise
about why. The reference does not state the criterion behind its sentence, but
it does state the quantities it reports for its simulations: the concentration
map $c(x,t)$, the activation time, which it defines through the threshold
$c_{\mathrm{crit}} = 0.95$, and the concentration integrated over the domain,
$C(t)=\int_{\mathcal{B}} c \,\mathrm{d}V$. The front position is not among
them. Our diagnostic is therefore our own choice and not theirs, the two
statements are about different quantities, and they cannot be compared directly.
We adopt the front position because the claim being reproduced concerns
spreading, and the location of the $c=0.5$ level measures directly how far the
solution has spread; what is reported here is the behaviour of that quantity and
nothing beyond it.

Panel (b) quantifies the same experiment with three dimensionless differences
from the finest run, all evaluated on the profiles at the common final time:
\begin{equation}
e_\infty(\Delta t) = \max_{x\in\Omega} \bigl| c_{\Delta t} - c_{\mathrm{ref}} \bigr|,
\qquad
e_2(\Delta t) = \left( \frac{1}{|\Omega|}
  \int_\Omega \bigl( c_{\Delta t} - c_{\mathrm{ref}} \bigr)^2 \,\mathrm{d}x
  \right)^{1/2},
\qquad
e_{\mathrm{f}}(\Delta t) = \bigl| x_{\mathrm{f}}(\Delta t)
  - x_{\mathrm{f}}(\mathrm{ref}) \bigr|,
\label{eq:timestepdiff}
\end{equation}
with $\Omega=(-1,1)$ and $c_{\mathrm{ref}}$ the profile of the run at
$\Delta t = 0.025$. The second is a root mean square rather than a plain $L^2$
norm, so that it is dimensionless and directly comparable with the first; the
third is a displacement in a domain of half-width one, hence also
dimensionless. The front difference $e_{\mathrm{f}}$ is defined only where a
crossing exists, so its curve stops at $\Delta t = 0.1$ and the shaded region
marks the steps for which the quantity does not exist at all. The two norms of
$c$ remain well defined for every step, and they saturate at the values they
take against the uniformly filled state.

Two things this experiment does not establish should be stated plainly. It is
not a study of the temporal order of convergence: the run at
$\Delta t = 0.025$ is a numerical reference and not an exact solution, so the
quantities in \eqref{eq:timestepdiff} are differences between two numerical
solutions and not errors. And the range of admissible steps is too narrow for
an order to be extracted from it, since only three of the six runs possess a
front and the remaining differences have saturated. The formal temporal orders
are verified separately in the paragraph below, against a problem with an exact
solution.

\begin{figure}[htbp]
\centering
\includegraphics[width=\textwidth]{images/time_step_study.pdf}
\caption{Time-step sensitivity at $\alpha=2$, $d=2\cdot10^{-4}$, with all six
runs compared at the common time $T=19.2$. (a) the six final profiles, with the
$c=0.5$ level marked by the dashed line and the resulting front positions by
the circles. Only three runs cross that level:
$x_{\mathrm{f}} = 0.6889$, $0.7268$ and $0.8145$ for
$\Delta t = 0.025$, $0.05$ and $0.1$. For $\Delta t \geq 0.2$ the interval is
full and no crossing exists, the minimum concentration being $0.9320$,
$0.99996$ and $0.9999999$ respectively; the profile at $\Delta t = 0.4$ is
dashed because it coincides with the one at $\Delta t = 0.3$ to $4\cdot10^{-5}$.
On the plateau $|x|\leq 0.5$ all six runs agree to $3.3\cdot10^{-4}$, so the
disagreement is confined to the front. (b) the three differences from the
finest run defined in \eqref{eq:timestepdiff}; the run at $\Delta t = 0.025$ is
the reference and has difference identically zero, so it carries no marker on
the logarithmic axis. The front difference $e_{\mathrm{f}}$ stops at
$\Delta t = 0.1$ because no front exists beyond it, as the shaded region
records; the two norms of $c$ are defined throughout. The finest run is a
numerical reference and not an exact solution, so no quantity shown is an error
and no convergence order is implied.}
\label{fig:timestep}
\end{figure}

Before moving to the connectome we verify both nonlinear time integrators
against an exact solution. On a spatially uniform datum with a uniform
coefficient the diffusion term vanishes and the exact solution is the scalar
logistic curve. Temporal refinement gives the expected orders: $1.99918$ and
$1.99966$ for the semi-implicit Crank--Nicolson scheme
of~\cite{corti2023}, and $1.00144$ and $1.00036$ for Backward Euler with
Newton.


% ---------------------------------------------------------------------
%  DA LEGGERE prima di 4.2
%    [F] par. 3.1 "The connectivity-weighted graph", pp. 4-5: la
%        costruzione del grafo dal Budapest Reference Connectome v3.0,
%        la mappatura da 1015 a 83 nodi, e TUTTE le statistiche che
%        riproduciamo nella Tabella~\ref{tab:connectome}. E' il paragrafo
%        piu' importante di tutto il capitolo per la parte di dati.
%    [F] par. 3.2 "The graph Laplacian", p. 5: equazione (3.1), la
%        definizione di A_IJ = n_IJ / l_IJ e del Laplaciano pesato.
%        Da confrontare con la nostra scelta di peso, discussa nel testo.
%    [C] par. 4.1 "Graph reconstruction and PET projection", p. 8: come
%        Corti costruisce il suo grafo da DWI e trattografia. Serve per
%        capire perche' NON possiamo riprodurre il loro grafo a 246 nodi.
%    Facoltativo ma utile: [C] par. 2.1, p. 3, prime righe, dove il grafo
%    e' introdotto come discretizzazione del dominio.
% ---------------------------------------------------------------------

\subsection{The brain connectome as a metric graph}
\label{sec:connectome}

Having verified the nonlinear solver on a controlled one-dimensional problem,
we move to the network on which the disease actually spreads.
Following~\cite{fornari2019}, we represent the brain as a weighted undirected
graph $\mathcal{G}=(\mathcal{V},\mathcal{E})$ whose vertices are anatomical
regions and whose edges are the axonal bundles connecting them. The graph is
extracted from the public Budapest Reference Connectome v3.0, which aggregates
the tractography of $418$ healthy subjects of the Human Connectome Project. We
retain the connections occurring in at least five subjects and aggregate the
$1015$ fine parcels into the $83$ regions of the FreeSurfer parcellation,
removing the connections internal to a parcel.

The reconstruction is verified against the published values before any
simulation is run. Table~\ref{tab:connectome} collects the comparison. The
topology is recovered exactly and the weighted quantities agree to within a few
percent. The residual discrepancy is expected and its origin is known: the
downloadable file provides the median consensus over the subjects, whereas the
published statistics are averages over the $418$ brains, and the original
$83\times83$ averaged matrix is not distributed. We also downloaded and
aggregated the mean-mode graph with the same confidence threshold; it gives a
mean adjacency of $2.4398$ and a mean weighted degree of $66.4339$, both
farther from the published values than the median reconstruction, so the median
graph is retained.

\begin{table}[htbp]
\centering
\begin{tabular}{lrr}
\hline
Quantity & This work & \cite{fornari2019} \\
\hline
fine graph vertices / edges         & $1015$ / $37477$    & $1015$ / $37477$ \\
aggregated vertices / edges         & $83$ / $1130$       & $83$ / $1130$ \\
degree range                        & $6$ -- $48$         & $6$ -- $48$ \\
mean fibre number                   & $40.16$             & $40.2$ \\
mean fibre length                   & $38.40$ mm          & $38.40$ mm \\
mean adjacency $\bar{A}_{IJ}$       & $1.6669$            & $1.57$ \\
adjacency range                     & $0.0092$ -- $36.87$ & $0.01$ -- $35.32$ \\
mean weighted degree $\bar{D}_{II}$ & $45.39$             & $42.8$ \\
weighted degree range               & $2.29$ -- $134.84$  & $2.1$ -- $127.6$ \\
\hline
\end{tabular}
\caption{Verification of the reconstructed $83$-region connectome against the
values published in~\cite{fornari2019}. The topology is exact; the weights are
close but not identical because the public file is a median consensus while the
reference statistics are averages over $418$ subjects.}
\label{tab:connectome}
\end{table}

One property of the weighting deserves attention here, because it reappears in
the regional results. The weight of an edge is the fibre count divided by the
fibre length, so it penalises long-range bundles. Summed over pairs of
anatomical groups, the connectivity between the frontal and the parietal group
reaches $183.5$ and between the parietal and the occipital group $201.8$,
while the frontal-to-temporal entry is $0.88$, distributed over $14$
connections none of which exceeds $0.19$. The short association fibres dominate the weighted
graph.

% ---------------------------------------------------------------------
%  ANALOGA di Figura~\ref{fig:regions}
%    [C] Fig. 2, p. 8, pannelli (c) e (d): i nodi del grafo colorati per
%        regione, e il grafo con gli archi. I pannelli (a) superficie MRI
%        e (b) trattografia NON sono riprodotti e il motivo va detto nel
%        testo: non abbiamo quei dati e il grafo non e' ricostruito da
%        immagini di paziente.
%    [F] Fig. 5, p. 5: e' l'altra figura da guardare, per lo stile del
%        rendering (grafo dentro una superficie piale traslucida, tre
%        orientamenti anatomici). La nostra figura prende la struttura da
%        [C] e il rendering da [F].
% ---------------------------------------------------------------------

Figure~\ref{fig:regions} shows the resulting domain, with the vertices at their
anatomical coordinates coloured by the seven anatomical groups used later in
this chapter, in the three anatomical projections.

\begin{figure}[htbp]
\centering
\includegraphics[width=\textwidth]{images/connectome_regions.pdf}
\caption{The $83$-region connectome, seen through the pial surface in the
sagittal, coronal and axial projections. Top: vertices coloured by anatomical
group. Bottom: connections coloured by connectivity weight. The surface is the
same file used for the ParaView output and shares the coordinate frame of the
vertices; nothing is rescaled to make the two agree. The MRI and tractography
panels of the corresponding figure in~\cite{corti2023} are not reproduced,
since neither dataset is available and the graph is taken from a public
consensus connectome rather than reconstructed from patient images.}
\label{fig:regions}
\end{figure}

% ---------------------------------------------------------------------
%  ANALOGA di Figura~\ref{fig:connectogram}
%    [C] Fig. 3, p. 9: i grafi locali delle sette regioni e il
%        connettogramma con la soglia al 5% del peso massimo. Noi usiamo
%        la stessa soglia del 5%, ed e' un dettaglio che vale la pena
%        citare nel testo perche' rende il confronto diretto.
%    Nota: [C] disegna anche i sette sottografi regionali separatamente;
%    noi li sostituiamo con la matrice 7x7 aggregata, che porta la stessa
%    informazione in forma quantitativa.
% ---------------------------------------------------------------------

Figure~\ref{fig:connectogram} shows the connectivity structure in the form used
in~\cite{corti2023}, with the same threshold of $5\%$ of the strongest
connection, together with the connectivity summed over each pair of anatomical
groups.

\begin{figure}[htbp]
\centering
\includegraphics[width=\textwidth]{images/connectome_connectogram.pdf}
\caption{Connectivity of the reconstructed graph. Left: the connections above
$5\%$ of the strongest one, on a circular layout grouped by anatomical group,
following the threshold used in~\cite{corti2023}. Right: connectivity summed
over each pair of groups.}
\label{fig:connectogram}
\end{figure}

% ---------------------------------------------------------------------
%  ANALOGA di Figura~\ref{fig:topology}
%    [F] Fig. 6, p. 6. Corrispondenza pannello per pannello:
%      nostro (a) <-> loro (a), grado non pesato;
%      nostro (b) <-> loro (b), grado pesato dalla connettivita';
%      nostro (c) <-> loro (c), matrice di adiacenza.
%    Da leggere insieme alla figura: [F] par. 3.2, p. 5, che commenta i
%    range (2.1 <= D_II <= 127.6, media 42.8; 0.01 <= A_IJ <= 35.32,
%    media 1.57) e la struttura small-world a blocchi. Sono esattamente i
%    numeri della nostra Tabella~\ref{tab:connectome}.
%    Differenze da dichiarare: loro usano una scala lineare tipo jet e
%    disegnano lo zero in blu; noi usiamo una scala logaritmica e mettiamo
%    lo zero fuori scala. Il motivo e' misurato ed e' scritto nel testo.
% ---------------------------------------------------------------------

Figure~\ref{fig:topology} reproduces the three connectivity diagnostics
of~\cite{fornari2019}: the non-weighted degree, the connectivity-weighted
degree, and the weighted adjacency matrix. The ranges and the block structure
are consistent with the small-world organisation described there, with strongly
connected hubs inside each hemisphere and comparatively few interhemispheric
connections.

Two conventions of our panel differ from the published one and are stated
because they change how the figure is read. The adjacency uses a base-$10$
logarithmic colour scale over the full non-zero range, with nothing clipped:
the $2260$ non-zero entries span $3.603$ decades and $77.1\%$ of them fall
below $5\%$ of the maximum, so on a linear ramp three quarters of the
connectome would render as a single indistinguishable colour. The price is that
the top of the range is compressed, which is why the bar carries decade ticks.
Absent connections are not on the scale at all: the $4629$ zero entries, the
$83$ diagonal ones among them, are drawn in one flat grey. Rows and columns
keep the solver's own node order, with no permutation or clustering, so that
the panel can be laid beside the published one and compared cell by cell.

\begin{figure}[htbp]
\centering
\includegraphics[width=\textwidth]{images/connectome_topology.pdf}
\caption{Connectivity diagnostics of the reconstructed graph, after
figure 6 of~\cite{fornari2019}. (a) non-weighted degree and (b)
connectivity-weighted degree, with sphere radius and colour encoding the same
quantity, each panel normalised to its own range. (c) weighted adjacency
$A_{IJ}$ on a base-$10$ logarithmic scale over the full non-zero range, with
absent connections drawn off-scale in grey and anatomical-group strips along
all four edges. Rows and columns are in the solver's node order: right
hemisphere $0$--$40$, left $41$--$81$, brainstem $82$.}
\label{fig:topology}
\end{figure}


% ---------------------------------------------------------------------
%  DA LEGGERE prima di 4.3
%    [F] par. 3.3 "The network model", p. 5: equazione (3.2), il modello
%        nodale con N incognite. E' il modello di riferimento con cui ci
%        confrontiamo, e va letto con attenzione perche' la differenza
%        con il nostro e' sottile: loro hanno una incognita per regione,
%        noi una funzione continua su tutto il grafo metrico.
%    [C] par. 2.1 "Discrete formulation", p. 3: equazione (3), che e' lo
%        stesso modello nodale, e il commento sul Laplaciano come
%        discretizzazione dell'operatore continuo.
%    [W] par. 2.3, p. 268, equazioni (10) e (11): la formulazione a
%        elementi finiti nel continuo tridimensionale. Serve per vedere
%        che il nostro termine di reazione integrato e' l'analogo su
%        grafo metrico del loro, e non un prodotto nodale.
%    Da confrontare esplicitamente: [F] eq. (3.2) contro la nostra
%    eq.~\eqref{eq:fem}. La differenza e' la matrice di massa.
% ---------------------------------------------------------------------

\subsection{Network model and metric-graph formulation}
\label{sec:formulation}

On the connectome, \cite{fornari2019} solve a system of ordinary differential
equations with one unknown per region,
\begin{equation}
\frac{\mathrm{d}c_I}{\mathrm{d}t}
  = -\sum_{J=1}^{N} L_{IJ}\,c_J + \alpha\,c_I\,[1-c_I],
  \qquad I = 1,\dots,N,
\label{eq:nodal}
\end{equation}
where $L = D - A$ is the connectivity-weighted graph Laplacian. This is a nodal
model: the brain regions are points and the connections carry no internal
structure. The same equation is the starting point of~\cite{corti2023}.

The formulation developed in this work replaces the graph by a metric graph.
Every connection becomes an interval carrying its own diffusion coefficient,
the concentration is a continuous piecewise linear function on the whole
network, and on each edge $e$ it satisfies
\begin{equation}
\partial_t c - \partial_s\!\left(D_e\,\partial_s c\right)
  = \alpha(s)\,c\,(1-c),
\end{equation}
with continuity at the vertices and the weighted Kirchhoff flux balance
\[
\sum_{e \ni v} D_e\,\partial_{n_e} c_e(v,t) = 0
\]
imposed naturally by the weak formulation. The P1 semi-discretization is
\begin{equation}
M\,\dot{\mathbf{C}} + K_D\,\mathbf{C} = \mathbf{R}(\mathbf{C}),
\qquad
R_i(\mathbf{C}) = \int_{\Gamma} \alpha\,c_h\,(1-c_h)\,\varphi_i \,.
\label{eq:fem}
\end{equation}
The reaction vector and its Jacobian are integrated with three-point
Gauss--Legendre quadrature on every cell, which is exact for the degree-four
integrands generated by a P1 concentration and a P1 reaction coefficient. The
reaction is therefore a consistent finite element integral and not a nodal
Hadamard product, which is the operation appearing in equation (4)
of~\cite{corti2023}.

Two differences between \eqref{eq:nodal} and \eqref{eq:fem} account for
essentially all the discrepancy observed below. First, the mass matrix $M$ is
not the identity, so a value prescribed at a vertex has a metric support that
depends on how many connections meet there. Second, one element per connection
makes the vertex unknowns of \eqref{eq:fem} coincide in number with the
unknowns of \eqref{eq:nodal} but not in meaning, and refining the mesh inside
the connections adds degrees of freedom that the nodal model cannot represent
at all. It is worth noting for later that the nodal model is exactly the
mass-lumped counterpart of \eqref{eq:fem}.


% ---------------------------------------------------------------------
%  DA LEGGERE prima di 4.4
%    [F] par. 4 "Biomarker models", p. 6, prime righe: la definizione di
%        biomarker abnormality.
%    [F] par. 4.1 "The Fisher-Kolmogorov model", p. 6: e' il paragrafo da
%        riprodurre. Contiene alpha = 0.5, il seed entorinale c0 = 0.1,
%        lo schema implicito con 100 passi da Delta t = 0.4 anni, e
%        l'equazione (4.1) del biomarcatore C(t) = somma di c_I.
%        Contiene anche l'affermazione sulla sequenza di attivazione
%        temporale -> frontale -> parietale -> occipitale, che e'
%        l'unica cosa che NON riproduciamo.
%  ANALOGA di Figura~\ref{fig:biomarker}
%    [F] Fig. 7, p. 7. Stessi quattro lobi, stessi colori (verde
%        temporale, rosso frontale, arancio parietale, blu occipitale),
%        stessa curva tratteggiata per la media di rete. La differenza
%        sostanziale, cioe' che le nostre quattro curve coincidono, e'
%        proprio il confronto da fare guardando le due figure affiancate.
%    Nota: i pannelli inferiori della nostra figura (scarto dalla media)
%    non hanno corrispondente in [F]; li abbiamo aggiunti perche' senza
%    di essi le curve sovrapposte non direbbero nulla.
% ---------------------------------------------------------------------

\subsection{Comparison with the network model}
\label{sec:comparison}

We reproduce the baseline simulation of~\cite{fornari2019}: conversion rate
$\alpha=0.5$, misfolded protein seeded in the entorhinal cortex at $c_0=0.1$,
all other regions healthy, and $100$ implicit steps of $\Delta t = 0.4$ years.
The biomarker abnormality of a group of regions is the concentration averaged
over its vertices, expressed as a percentage.

Both solutions remain in the physical interval, up to a transient undershoot of
order $10^{-4}$ in the finite element solution which we trace below to the
coarse mesh, and both approach the stable state $c=1$. The biomarker curves
have the expected sigmoid shape, and the $50$-percent crossing of the network
average occurs at $11.11$ years for the nodal model and $12.68$ years for the
finite element model, the difference being the metric-support effect described
in Section~\ref{sec:formulation}.

The activation sequence, however, is not reproduced. \cite{fornari2019} report
a separation of roughly five years between the temporal and the occipital lobe;
our four lobe curves cross the $50$-percent level within $2.9\cdot10^{-7}$
years of each other in the nodal model and $5.7\cdot10^{-3}$ years in the
finite element model. Since our nodal model is \eqref{eq:nodal} implemented
literally, the discrepancy cannot be a defect of the spatial discretization,
and Section~\ref{sec:damkohler} is devoted to explaining it.

\begin{figure}[htbp]
\centering
\includegraphics[width=\textwidth]{images/biomarker_comparison.pdf}
\caption{Biomarker abnormality of the four lobes and of the network, for the
nodal reference and for the metric-graph finite element model, at $\alpha=0.5$
with $100$ steps of $\Delta t = 0.4$ years. Since the four regional curves
nearly coincide, the lower panels report each curve minus the network average;
no rescaling is applied to the upper panels.}
\label{fig:biomarker}
\end{figure}

% ---------------------------------------------------------------------
%  DA LEGGERE prima di 4.4.1
%    Nessuna lettura nei tre paper: nessuno dei tre riporta uno studio di
%    raffinamento in spazio o in tempo su connettoma. E' interamente
%    nostro, ed e' esattamente il tipo di contributo che ci si aspetta da
%    un lavoro di metodi numerici applicato a un modello preso dalla
%    letteratura. Vale la pena dirlo nel testo.
%    Se serve un termine di paragone sul ruolo di Delta t, l'unico e'
%    [W] par. 2.4, p. 268, gia' letto per la Sezione 4.1.
% ---------------------------------------------------------------------

\subsubsection{Convergence}
\label{sec:convergence}

None of the reference works reports a refinement study on the connectome, so we
carry one out before drawing any conclusion from the comparison above. The
spatial study keeps $\Delta t = 0.4$ and refines each connection into $1$, $2$,
$4$ and $8$ elements. The maximum deviation of the biomarker curves from the
finest run decreases as
\[
0.4389\%, \quad 0.1178\%, \quad 0.0229\%, \quad 0,
\]
over $83$, $1213$, $3473$ and $7993$ degrees of freedom. The small negative
transient disappears from two elements per connection onwards, confirming that
it is a coarse-mesh artefact and not a failure of the nonlinear solver.

The temporal study uses eight elements per connection and the steps
$\Delta t = 0.8, 0.4, 0.2, 0.1, 0.05$; successive differences between biomarker
curves give the observed rates $1.0552$, $1.0263$ and $1.0126$, converging to
the first order expected of the Backward Euler scheme.

We retain $\Delta t = 0.4$ in the comparison above in order to match the
discretization of the reference, but the study just described moves the network
crossing time from $11.50$ years at $\Delta t = 0.8$ to $13.70$ years at
$\Delta t = 0.05$. The absolute crossing times therefore carry a
time-discretization error of order one year, in our simulation as well as in
the reference, which uses the same step with a first-order scheme. Agreement of
absolute times between the two works should be read with that in mind.


% ---------------------------------------------------------------------
%  DA LEGGERE prima di 4.5
%    [F] par. 3.3, p. 5, di nuovo: e' l'equazione a cui aggiungiamo il
%        fattore di scala rho. Rileggerla con in mente che nel paper
%        rho = 1 implicitamente.
%    [F] par. 4.1, p. 6: la sequenza di attivazione dichiarata, che e'
%        cio' che questa sezione spiega.
%    Nessun altro riferimento: l'analisi con il numero di Damkohler e'
%    nostra. Non esiste in nessuno dei tre paper, ed e' il contributo
%    principale del capitolo. Conviene presentarla come tale.
%  ANALOGA di Figura~\ref{fig:damkohler}
%    NESSUNA. E' una figura originale. I pannelli superiori usano la
%    stessa convenzione grafica di [F] Fig. 7, p. 7, cosi' che il
%    lettore possa confrontarli, ma il pannello inferiore non ha
%    corrispondente in letteratura.
% ---------------------------------------------------------------------

\subsection{What controls the regional separation}
\label{sec:damkohler}

Consider \eqref{eq:nodal} with the Laplacian scaled by a positive factor
$\rho$, so that $\rho=1$ recovers the published model. Linearising around a
small uniform concentration, the average over the network grows at the reaction
rate $\alpha$, while any imbalance between regions relaxes at the rate
$\rho\lambda_2$, where $\lambda_2$ is the Fiedler value of $L$, its smallest
non-zero eigenvalue. A difference between lobes can therefore survive long
enough to be visible only if the reaction is faster than the homogenisation of
the graph, that is if the Damk\"ohler number
\begin{equation}
\mathrm{Da} = \frac{\alpha}{\rho\,\lambda_2}
\label{eq:damkohler}
\end{equation}
is large. For the reconstructed connectome we compute $\lambda_2 = 0.7954$ and
$\lambda_{\max} = 154.04$.

At the parameters of the reference, $\rho = 1$ and $\alpha = 0.5$, we obtain
$\mathrm{Da} = 0.63$. The connectome homogenises faster than the reaction
develops, and coincident lobe curves are the correct solution of
\eqref{eq:nodal} at unit weight scale, not a numerical artefact.
Figure~\ref{fig:damkohler} quantifies this by sweeping $\rho$: the lobe
separation grows from $2.9\cdot10^{-7}$ years at $\mathrm{Da}=0.63$ to $6.92$
years at $\mathrm{Da}=629$, seven decades of variation with the topology, the
seed and the conversion rate held fixed. Recovering the separation reported
in~\cite{fornari2019} requires $\mathrm{Da}$ of order $10^{2}$, that is $\rho$
close to $0.005$. We conclude that the published figure cannot be produced by
the connectivity-weighted Laplacian used at unit scale, whatever reconstruction
of the weights is adopted.

\begin{figure}[htbp]
\centering
\includegraphics[width=\textwidth]{images/diffusion_scaling.pdf}
\caption{Regional separation as a function of the Damk\"ohler number
\eqref{eq:damkohler}. Top: lobe biomarkers at three scalings, drawn with the
convention of figure 7 of~\cite{fornari2019} so that the two can be compared.
Bottom: the separation between the earliest and the latest lobe against
$\mathrm{Da}$, with the two settings used in this chapter marked and the shaded
region indicating where the consistent-mass P1 discretization loses
boundedness.}
\label{fig:damkohler}
\end{figure}

This also reconciles two results of this chapter that would otherwise look
contradictory. The deterministic model of Section~\ref{sec:corti} normalises
the diffusivity as $D_e = w_e/\max_e w_e$, that is $\rho = 0.0271$, and uses
regional reaction rates whose average over the $83$ vertices is $0.1252$; its
Damk\"ohler number is $5.80$, on the other side of the crossover, which is why
its regional curves separate clearly while those of
Section~\ref{sec:comparison} do not.

\subsubsection{A validity boundary for the discretization}
\label{sec:boundary}

The same sweep bounds the finite element formulation. Above $\mathrm{Da}$ of
about $13$, that is $\rho \le 0.02$ at $\alpha=0.5$, the discrete concentration
leaves $[0,1]$ and then diverges: at $\rho = 0.005$ the nodal reference stays
within $[0.99987, 1]$ while the finite element solution reaches $[-922, 308]$.
The failure occurs identically with Backward Euler and with the semi-implicit
Crank--Nicolson scheme, so it is not a failure of the Newton solver, and it
persists at $1$, $2$, $4$ and $8$ elements per connection, so it is not a
matter of resolving the front.

The cause is the consistent mass matrix. Its off-diagonal entries are positive,
so the semi-discrete system does not satisfy a discrete maximum principle;
while diffusion dominates, the diffusion matrix compensates, and once the
reaction dominates nothing does. As observed in Section~\ref{sec:formulation},
the nodal model \eqref{eq:nodal} and the Hadamard product used
in~\cite{corti2023} are exactly the mass-lumped counterparts of
\eqref{eq:fem}, which is why they remain bounded at every scaling tested.
The same defect is visible on the one-dimensional problem of
Section~\ref{sec:onedim}, which is worth recording because it shows that the
boundary is a property of the discretization and not of the connectome: with
the sharp single-node initial datum, refining the time step at $\alpha=2$ and
$d=10^{-4}$ makes the undershoot grow rather than shrink, from
$-1.5\cdot10^{-13}$ at $\Delta t = 0.1$ to $-5.4\cdot10^{-6}$ at $0.05$ and
$-9.1\cdot10^{-4}$ at $0.025$. Introducing an optional lumped mass matrix is
therefore the natural way to extend the present formulation into the
reaction-dominated regime, and we leave it as the immediate continuation of
this work. Both simulations reported in this chapter
satisfy $\mathrm{Da}\le 6$ and are therefore well inside the region where the
discretization is reliable.


% ---------------------------------------------------------------------
%  DA LEGGERE prima di 4.6
%    [F] par. 5 "Treatment opportunities", p. 7, introduzione.
%    [F] par. 5.1 "Delaying conversion", pp. 7-8: e' l'esperimento che
%        riproduciamo. Contiene l'interpretazione di alpha come
%        combinazione di k12, k0, k1 e k1-tilde tramite l'equazione (2.7)
%        di p. 3, che vale la pena citare nel testo per dare senso
%        fisico alla variazione di alpha.
%        Contiene anche la conclusione che riproduciamo alla lettera:
%        abbassare alpha ritarda ma non crea uno stato stabile
%        intermedio.
%  ANALOGA di Figura~\ref{fig:alpha}
%    [F] Fig. 11, p. 8. Stesso insieme di valori alpha = 0, 0.1, ..., 0.5
%        e stesso orizzonte di 80 anni. La nostra aggiunge il confronto
%        nodale/FEM, che in [F] non esiste perche' hanno un solo modello.
%    Da guardare anche [F] Fig. 10, p. 8, che e' la versione anatomica
%    dello stesso esperimento; la nostra corrispondente e' la
%    Figura~\ref{fig:progression}, discussa piu' avanti.
% ---------------------------------------------------------------------

\subsection{Sensitivity to the conversion rate}
\label{sec:alpha}

We repeat the conversion-rate experiment of~\cite{fornari2019} with
$\alpha = 0, 0.1, \dots, 0.5$ over $80$ years. In the reference the conversion
rate is not a free parameter but a combination of production, clearance and
misfolding rates, so lowering it is the model's representation of a therapeutic
intervention.

The qualitative conclusion of the reference is recovered: lowering $\alpha$
delays the transition but does not create an intermediate stable state, and
every positive conversion rate drives the network to $c=1$. The $50$-percent
crossing of the network average moves from $11.11$ years at $\alpha = 0.5$ to
$59.33$ years at $\alpha = 0.1$ in the nodal model, and from $12.68$ to $67.81$
years in the finite element model. For $\alpha = 0$ the reaction vanishes and
the two conserved quantities are preserved to round-off: the nodal average
drifts by $3.9\cdot10^{-16}$ and the normalised finite element metric integral
$\mathbf{1}^{\!\top} M \mathbf{C}/\mathbf{1}^{\!\top} M \mathbf{1}$ by
$9.7\cdot10^{-17}$.

The nodal and finite element curves start from different global averages, for
the reason given in Section~\ref{sec:formulation}. No amplitude rescaling is
applied to hide this difference, since it is a property of the two models and
not an error of either.

\begin{figure}[htbp]
\centering
\includegraphics[width=\textwidth]{images/alpha_sensitivity.pdf}
\caption{Conversion-rate sensitivity, after figure 11 of~\cite{fornari2019},
for the nodal reference and the metric-graph finite element model over $80$
years. The two models start from different global averages because a value
prescribed at a vertex has a connectivity-dependent metric support in the
finite element domain; no rescaling is applied.}
\label{fig:alpha}
\end{figure}


% ---------------------------------------------------------------------
%  DA LEGGERE prima di 4.7
%    [F] par. 4.3 "Infection times", p. 7: e' l'esperimento che
%        riproduciamo. Attenzione a due cose che vanno dette nel testo:
%        (i) loro lo fanno con il modello Heterodimer, non con
%        Fisher-Kolmogorov, quindi i tempi assoluti non sono
%        confrontabili e noi confrontiamo solo l'ordinamento;
%        (ii) la frase sulla corteccia entorinale che ha il secondo tempo
%        di infezione piu' lungo, e l'argomento clinico che ne traggono,
%        e' esattamente il risultato che ritroviamo indipendentemente.
%        Vale la pena citarla per esteso.
%  ANALOGA di Figura~\ref{fig:vulnerability}
%    [F] Fig. 9, p. 8. Loro mostrano le 83 curve colorate per tempo di
%        infezione e le sfere colorate sull'anatomia. I nostri pannelli
%        (a) e (b) sono gli analoghi diretti.
%    I nostri pannelli (c) e (d) non hanno corrispondente: (c) e' il
%    meccanismo (tempo contro grado pesato) e (d) e' il confronto fra i
%    due scaling. Sono aggiunte nostre e vanno presentate come tali.
% ---------------------------------------------------------------------

\subsection{Regional vulnerability}
\label{sec:vulnerability}

The network enters the model only through $L$, so it is natural to ask which
regions are the most dangerous places for the disease to start. Following the
infection-time study of~\cite{fornari2019} we seed each of the $83$ regions in
turn and record the time at which the network biomarker reaches $95\%$. The
reference performs this experiment with the Heterodimer model rather than with
Fisher--Kolmogorov, so the absolute times are not comparable and we compare only
the ranking, which depends on the network rather than on the kinetics.

At $\rho = 1$ the infection times span $0.02$ years, which is the homogenisation
effect quantified in Section~\ref{sec:damkohler}; at $\rho=0.02$ they span four
years and the ordering becomes measurable. Three observations of the reference
survive at both scalings.

The fastest seeds are subcortical hubs: we obtain the two caudate nuclei
followed by the two thalami, where the reference reports the putamen and the
insula. Both identify the same feature of the network, and the reference itself
notes that the caudate carries the largest number of connections in the graph.
The slowest seeds are the frontal and temporal poles; the frontal pole has the
smallest weighted degree of the entire graph, $2.29$. Finally, and most
interestingly, the entorhinal cortex, the region where the pathology is first
observed clinically and the seed used throughout this chapter, ranks $76$th and
$80$th of $83$ at $\rho=1$ and $78$th and $81$st at $\rho=0.02$: it is one of
the least efficient possible starting points. The reference draws from the same
observation the argument that this is part of why the early stages of the
disease are so hard to detect, and it is recovered here from an independent
reconstruction of the graph.

The third panel of the figure gives the mechanism. The infection time decreases
monotonically with the weighted degree of the seeded vertex: a poorly connected
seed keeps its concentration local, and the quadratic term of the reaction then
slows the growth of the network average.

\begin{figure}[htbp]
\centering
\includegraphics[width=\textwidth]{images/seeding_vulnerability.pdf}
\caption{Regional vulnerability, after figure 9 of~\cite{fornari2019}. Each of
the $83$ regions is seeded in turn. (a) the resulting network biomarker curves
and (b) the same quantity on the anatomy, both at $\rho=0.02$; (c) infection
time against the weighted degree of the seeded region; (d) the delay of every
seed behind the fastest one, at both scalings, on a logarithmic axis. The
anatomy is not shown at $\rho=1$: there the $83$ infection times differ by
$0.02$ years, and colouring the regions by that quantity would paint them all
the same.}
\label{fig:vulnerability}
\end{figure}


% ---------------------------------------------------------------------
%  DA LEGGERE prima di 4.8
%    [C] par. 2 "Mathematical model", pp. 2-3: equazioni (1) e (2), cioe'
%        il coefficiente di reazione come campo costante a tratti sulle
%        sette regioni.
%    [C] par. 2.1 "Discrete formulation", p. 3: equazione (4), lo schema
%        semi-implicito di Crank-Nicolson con estrapolazione del secondo
%        ordine. E' lo schema che implementiamo, e va confrontato riga
%        per riga con la nostra eq.~\eqref{eq:corti}.
%    [C] Tabella 1, p. 11: le sette medie stimate mu_l. Sono i numeri che
%        usiamo. Da citare esplicitamente.
%    [C] par. 4.1, p. 8: la proiezione PET che costituisce il dato
%        iniziale, e che noi NON abbiamo. Leggere per capire cosa
%        stiamo sostituendo con il dato sintetico.
%    [C] par. 4.3 "Prediction of disease development", p. 11: l'orizzonte
%        di 20 anni e dt = 0.2 (loro scrivono 2e-1), e dt = 0.02 per
%        l'analisi forward stocastica.
%    [C] Tabella 3, p. 13: le concentrazioni attese per lobo a
%        t = 5, 10, 15, 20. E' la tabella con cui confrontiamo la nostra
%        classifica regionale.
%    Attenzione: la prosa di [C] a p. 12 dice "higher concentration in
%    the parietal and frontal lobes", mentre la loro Tabella 3 mette il
%    parietale penultimo. Usiamo la tabella, che e' inequivocabile, e lo
%    diciamo.
% ---------------------------------------------------------------------

\subsection{Deterministic model with regional reaction rates}
\label{sec:corti}

The last application follows~\cite{corti2023}, who solve a patient-specific
stochastic Fisher--Kolmogorov model on a $246$-node Brainnetome graph and
calibrate the reaction coefficient region by region against two PET
acquisitions. Neither the patient graph nor the PET images are available to us,
so we reproduce the mathematical model, the time discretization and the
estimated mean parameters, and we state clearly that we do not reproduce the
clinical trajectory of their patient.

The reaction coefficient is piecewise constant over seven anatomical groups and
takes the mean values estimated in table 1 of the reference: $0.1801$ in the
frontal lobe, $0.1421$ temporal, $0.0627$ parietal, $0.1005$ insular, $0.1351$
limbic, $0.0545$ occipital and $0.1147$ in the subcortical nuclei. The nodal
values are interpolated linearly along each connection. The diffusivity is the
normalised connectivity $D_e = w_e/\max_e w_e$; the reference does not state a
normalisation, and Section~\ref{sec:damkohler} shows that this choice is what
places the simulation in the regime where regional curves separate.

The time discretization is the semi-implicit Crank--Nicolson scheme of the
reference: diffusion and the first concentration factor are centred in time,
and only the saturation factor is extrapolated,
\begin{equation}
\mathbf{C}^\star = \tfrac{3}{2}\mathbf{C}^{n} - \tfrac{1}{2}\mathbf{C}^{n-1},
\qquad
G(\mathbf{C}^\star)_{ij}
  = \int_\Gamma \alpha\,(1-c^\star)\,\varphi_i\varphi_j ,
\label{eq:corti}
\end{equation}
\begin{equation}
\Big(M + \tfrac{\Delta t}{2} K_D - \tfrac{\Delta t}{2} G(\mathbf{C}^\star)\Big)
  \mathbf{C}^{n+1}
= \Big(M - \tfrac{\Delta t}{2} K_D + \tfrac{\Delta t}{2} G(\mathbf{C}^\star)\Big)
  \mathbf{C}^{n},
\end{equation}
started from $\mathbf{C}^{-1} = \mathbf{C}^{0}$. Since $G$ changes at every
step, the sparse time-step matrix is rebuilt and factorised at every step, a
point we return to in Section~\ref{sec:performance}. We use $\Delta t = 0.2$
years, four elements per connection, $3473$ degrees of freedom and a horizon of
$20$ years, and we replace the unavailable PET projection by a documented
synthetic datum: entorhinal and hippocampal regions start at $c = 0.10$ and all
others at $c = 0.01$.

The solution stays in $[0.0250, 0.6536]$ throughout, every regional average
grows monotonically, and the mean concentration sampled at the $83$ anatomical
vertices rises from $0.0143$ to $0.1578$. The absolute levels are far below
those of the reference, as they must be given the different initial datum, but
the ordering of the regions can be compared, and Table~\ref{tab:corti} does so.
Twenty of the twenty-one pairwise orderings agree, a rank correlation of
$0.96$. The single inversion is frontal against limbic and is explained by the
initial condition: the synthetic seed places $c=0.10$ on entorhinal and
hippocampal vertices, all of which belong to the limbic group. Everything below
the leading pair is reproduced in the published order using nothing but the
seven reaction coefficients and the connectome topology.

\begin{table}[htbp]
\centering
\begin{tabular}{lrrrr}
\hline
Group & \cite{corti2023}, $t=20$ & rank & this work & rank \\
\hline
frontal      & $0.9289$ & $1$ & $0.2021$ & $2$ \\
limbic       & $0.8905$ & $2$ & $0.2101$ & $1$ \\
temporal     & $0.8699$ & $3$ & $0.1670$ & $3$ \\
insular      & $0.8558$ & $4$ & $0.1292$ & $4$ \\
subcortical  & $0.8413$ & $5$ & $0.1270$ & $5$ \\
parietal     & $0.7738$ & $6$ & $0.0967$ & $6$ \\
occipital    & $0.7336$ & $7$ & $0.0735$ & $7$ \\
\hline
\end{tabular}
\caption{Regional concentrations at the final time, against table 3
of~\cite{corti2023}. The absolute values are not comparable, since the
reference starts from a patient PET projection and we start from a synthetic
datum, but the ordering is: twenty of the twenty-one pairwise orderings agree.
Note that the prose of the reference describes higher concentrations in the
parietal and frontal lobes, whereas its own table places the parietal group
second from last; we compare against the table.}
\label{tab:corti}
\end{table}

% ---------------------------------------------------------------------
%  ANALOGA di Figura~\ref{fig:regional}
%    [C] Fig. 7, p. 14: le curve per lobo. Loro mostrano media e banda di
%        deviazione standard perche' il loro modello e' stocastico; noi
%        abbiamo un modello deterministico e mostriamo la sola media, che
%        e' l'analogo della loro curva centrale.
%    Da leggere insieme: [C] Tabella 3, p. 13, gia' citata sopra.
% ---------------------------------------------------------------------

\begin{figure}[htbp]
\centering
\includegraphics[width=0.85\textwidth]{images/regional_averages.pdf}
\caption{Regional averages of the deterministic model over $20$ years. The
reference reports mean and standard-deviation bands because its reaction
coefficient is a random field; the model solved here is the deterministic
counterpart obtained by fixing each coefficient at its estimated mean, so only
the central curve has a counterpart.}
\label{fig:regional}
\end{figure}


% ---------------------------------------------------------------------
%  DA LEGGERE prima di 4.9
%    [W] par. 3.3 "Prion-like spreading in neurodegenerative disease",
%        p. 273: descrive le simulazioni tridimensionali di cui le
%        Figure 10-12 sono il risultato.
%    [W] Fig. 11, p. 276, e Fig. 12, p. 277: sono le due figure a cui la
%        nostra si ispira. Guardare in particolare come il mesh
%        tetraedrico resta visibile in grigio dove la concentrazione e'
%        bassa: e' il meccanismo che abbiamo copiato, non la palette.
%    [W] Fig. 10, p. 275: la versione bidimensionale su sezione
%        sagittale, ancora piu' vicina alla nostra come composizione.
%    [F] par. 5.1 e Fig. 10, p. 8: l'altra figura anatomica di
%        riferimento, con l'effetto di ritardo al calare di alpha.
%    Da tenere presente leggendo: le loro immagini sono sezioni di un
%    campo continuo su mesh di volume con circa 80 000 incognite. La
%    differenza con il nostro substrato va dichiarata, non attenuata.
% ---------------------------------------------------------------------

\subsection{Anatomical progression}
\label{sec:progression}

Figure~\ref{fig:progression} shows the spreading itself, as a mid-sagittal
section in the construction of~\cite{weickenmeier2019}: the pial surface is
clipped to one half-space and the graph to the other, so the network sits
inside an opaque brain rather than floating over a translucent one, and the
zero of the colour ramp is the neutral grey of the discretisation, so the
lesion emerges from a pale wireframe instead of the whole object being dark.
Nothing is thresholded away; the white marks are the exact $c=0.5$ level set of
the P1 interpolant. The three stages are the stored times at which the network
mean concentration reaches $0.2$, $0.4$ and $0.6$, which is the same biomarker
the reference plots, so the staging is a property of the solution rather than a
choice of ours.

What is drawn deserves emphasis: not the $83$ regions, but every degree of
freedom the metric-graph formulation carries, the region vertices together with
the nodes along each connection, $10\,170$ in all at eight elements per
connection. This is precisely what the nodal model \eqref{eq:nodal} does not
have, and it is what makes the field read as a volume rather than as a
wireframe.

\begin{figure}[htbp]
\centering
\includegraphics[width=\textwidth]{images/anatomical_progression.pdf}
\caption{Mid-sagittal section of the spreading, for the deterministic model of
Section~\ref{sec:corti} refined to eight elements per connection, in the
construction of figures 10 and 11 of~\cite{weickenmeier2019}. The pial surface
is clipped to $x\ge0$ and the graph to $x\le0$ on the same plane, so the
network is occluded by the anatomy. Every drawn element is a degree of freedom
of the simulation and tube and vertex radii are constant, so thickness encodes
nothing. The zero of the colour ramp is the grey of the discretisation, the
scale is fixed on $[0,1]$ in all panels, and nothing is thresholded away; white
marks are the exact $c=0.5$ level set. Stages are the stored times at which the
network mean reaches $0.2$, $0.4$ and $0.6$. The horizon is $60$ years rather
than the $20$ of Section~\ref{sec:corti}; model, coefficients, seed and time
step are unchanged.}
\label{fig:progression}
\end{figure}

% ---------------------------------------------------------------------
%  ANALOGA di Figura~\ref{fig:activation}
%    NESSUNA diretta. L'idea di una mappa dei tempi di attivazione compare
%    in [W] Fig. 3, p. 269, ma li' e' monodimensionale e nel piano (x,c);
%    la versione su rete e' nostra.
%    Il risultato che la figura porta, cioe' che l'ordine e' dettato dai
%    tassi di reazione e non dalla propagazione, non ha corrispondente in
%    nessuno dei tre paper ed e' il secondo contributo del capitolo dopo
%    il numero di Damkohler.
% ---------------------------------------------------------------------

Figure~\ref{fig:activation} condenses the whole history into one image by
colouring each degree of freedom with the year its concentration reaches
$0.5$, and it corrects an interpretation the progression invites. The obvious
reading of an anterior-to-posterior gradient is a front travelling outward from
the seed. That reading is wrong here. Over the $83$ region vertices the
activation time correlates with the local reaction coefficient at $r=-0.841$
and with the graph distance from the seed at $r=+0.072$, and the graph
eccentricity from the seed is $2$: every region lies within two connections of
the entorhinal and hippocampal seed. The order in which the network is
recruited is set by the regional conversion rates, $0.1801$ in the frontal
group against $0.0545$ in the occipital one, and not by propagation. This is
the conclusion of Section~\ref{sec:damkohler} reached independently, and it is
why no directional arrows are drawn: an arrow would assert a direction of
travel the solution does not contain.

\begin{figure}[htbp]
\centering
\includegraphics[width=0.86\textwidth]{images/activation_time.pdf}
\caption{Order in which the network is recruited: every degree of freedom
coloured by the first stored year at which its concentration reaches $0.5$. Ten
of the $10\,170$ degrees of freedom never reach it within the horizon and take
the late end of the ramp. The activation time correlates with the local
reaction coefficient at $r=-0.841$ and with the graph distance from the seed at
$r=+0.072$, so the anterior-to-posterior gradient is a reaction-rate pattern
and not a travelling front.}
\label{fig:activation}
\end{figure}

One limitation is structural and we state it rather than work around it. The
reference fills its brain silhouette because its unknown lives on a tetrahedral
volume mesh with about $80\,000$ unknowns; ours lives on $1130$
one-dimensional connections, so the coloured structure occupies a lens-shaped
core with bare tissue around it. That rim is an absence of degrees of freedom,
not an absence of disease. Filling it, by interpolating the graph field into
the volume or by assigning cortical triangles to their nearest region, is the
one operation that would make our figure resemble theirs, and it would
manufacture spatial detail the model does not carry. Two further points a
reader will otherwise misread: pale grey means low and never healthy, since the
initial condition is $c=0.01$ everywhere and the network minimum rises
monotonically, whereas the reference's white tissue genuinely is $c=0$; and the
interior nodes of a connection lie on the straight chord between two region
centroids, which is the discretization the solver used and not an anatomical
fibre path.


% ---------------------------------------------------------------------
%  DA LEGGERE prima di 4.10
%    [F] par. 4.1, p. 6, ultima frase: "runs 0.55 and 0.64 s without and
%        with output on a standard laptop computer". E' l'unico dato di
%        costo computazionale nei tre paper, ed e' il termine di
%        paragone che citiamo.
%    [F] par. 1, p. 2: il confronto con i modelli a elementi finiti
%        tridimensionali, 400 000 tetraedri e 80 000 gradi di liberta'.
%        Utile per inquadrare l'ordine di grandezza.
%    Nessun altro: ne' [C] ne' [W] riportano tempi di calcolo.
%  ANALOGA di Figura~\ref{fig:performance}
%    NESSUNA. Nessuno dei tre paper mostra misure di prestazione.
% ---------------------------------------------------------------------

\subsection{Sequential baseline}
\label{sec:performance}

The deterministic model above is also the reference problem for the performance
analysis that opens the next chapter. Timings are taken in a separate release
build with visualisation and diagnostic output disabled, thread counts pinned
to one, and reported as the median of five runs after one warm-up.

\begin{table}[htbp]
\centering
\begin{tabular}{rrrrr}
\hline
Elements per connection & Unknowns & Matrix nonzeros & $100$ steps [s] &
Per step [s] \\
\hline
$1$ & $83$   & $2343$  & $0.0914$ & $0.00091$ \\
$2$ & $1213$ & $5733$  & $1.7390$ & $0.01739$ \\
$4$ & $3473$ & $12513$ & $1.7381$ & $0.01738$ \\
$8$ & $7993$ & $26073$ & $2.4048$ & $0.02405$ \\
\hline
\end{tabular}
\caption{Sequential baseline for the deterministic model. Assembly accounts for
less than one percent of the total in every refined case.}
\label{tab:performance}
\end{table}

Two facts emerge from Table~\ref{tab:performance} and they set the agenda for
what follows. First, more than $99\%$ of the time is spent in the time loop and
almost none in assembly or input, so the target for optimisation is the
repeated sparse factorisation forced by the semi-implicit scheme
\eqref{eq:corti}, not the finite element assembly. Second, the cost per step is
not a function of the number of unknowns alone: the two-element mesh, with
$1213$ unknowns, costs exactly as much as the four-element mesh with $3473$.
The one-element mesh has a hub-dominated sparsity pattern inherited from the
connectome, and sparse factorisation depends on the elimination ordering and
the resulting fill-in at least as much as on the matrix dimension.

For scale, \cite{fornari2019} report that their nodal model runs its $100$
implicit steps in $0.55$ s without output on a standard laptop; our
metric-graph formulation at the same $83$ unknowns takes $0.091$ s for its
$100$ semi-implicit steps. The machines and the schemes differ, so this is an
order-of-magnitude check rather than a comparison, but it confirms that the P1
metric-graph formulation carries no cost penalty over the nodal network model
at the coarsest mesh, while allowing arbitrary refinement inside the
connections.

\begin{figure}[htbp]
\centering
\includegraphics[width=\textwidth]{images/sequential_performance.pdf}
\caption{Sequential baseline. Left: assembly, time loop and total against the
number of unknowns, with the band spanning the minimum and maximum of five
runs. Right: cost per time step, showing that it is not a function of the
unknown count alone.}
\label{fig:performance}
\end{figure}

This is precisely the observation that motivates the reordering study of the
next chapter, where bandwidth, profile, factor fill and factorisation time are
measured for different orderings before any parallel scaling result is
interpreted.
